\documentclass[conference]{IEEEtran}
\IEEEoverridecommandlockouts

%% ─── Packages ───────────────────────────────────────────────────────────────
\usepackage{cite}
\usepackage{amsmath,amssymb,amsfonts}
\usepackage{algorithmic}
\usepackage{graphicx}
\setkeys{Gin}{draft}
\usepackage{textcomp}
\usepackage{xcolor}
\usepackage{booktabs}
\usepackage{multirow}
\usepackage{array}
\usepackage{url}
\usepackage{microtype}
\usepackage{balance}
\usepackage{subfig}
\usepackage{float}
\usepackage{enumitem}
\usepackage{tabularx}
\usepackage{makecell}

\graphicspath{{figs/}}
\def\BibTeX{{\rm B\kern-.05em{\sc i\kern-.025em b}\kern-.08em T\kern-.1667em\lower.7ex\hbox{E}\kern-.125emX}}

%% ─── Custom colours ─────────────────────────────────────────────────────────
\definecolor{ieeeblu}{RGB}{26,82,118}
\definecolor{ieeeblk}{RGB}{30,30,30}

\begin{document}

%% ════════════════════════════════════════════════════════════════════════════
%%  TITLE
%% ════════════════════════════════════════════════════════════════════════════
\title{%
  Network Intrusion Detection Using Gradient Boosting and Deep\\
  Learning on the CIC-IDS Dataset: A Comprehensive Comparative Study
}

\author{%
  \IEEEauthorblockN{%
    Priyanshu Kumar Senapati~(1RUA24CSE0337),
    Rishita Bramhapuriya~(1RUA24CSE0368),
    Rao Kumar Srinidhi Reddy~(1RUA24CSE0361),
    Sanjana M Bhat~(1RUA24CSE0411)
  }
  \IEEEauthorblockA{%
    \textsuperscript{1}Department of Computer Science, University of Technology\\
    \textsuperscript{2}Department of Electrical Engineering, Institute of Technology\\
    \textsuperscript{3}Cybersecurity Research Lab, National Research Centre\\
    \{author1, author2, author3\}@university.edu
  }
}

\maketitle

%% ════════════════════════════════════════════════════════════════════════════
%%  ABSTRACT
%% ════════════════════════════════════════════════════════════════════════════
\begin{abstract}
This paper presents a rigorous comparative study of eight machine learning and
deep learning architectures for network intrusion detection and multi-class
attack classification on the CIC-IDS 2017--2019 benchmark, comprising
over~9.2~million labeled network flow records spanning eight traffic classes.
We evaluate two gradient boosting frameworks (LightGBM, XGBoost) and six
deep learning models (Multilayer Perceptron, TabNet, FT-Transformer-Lite,
1D-Convolutional Neural Network, Gated Recurrent Unit, and Autoencoder)
across two tasks: binary anomaly detection and fine-grained attack family
classification. LightGBM achieves the best binary AUC of~0.9973 and
multi-class Weighted~F\textsubscript{1} of~0.873, while TabNet leads among
deep learning models at AUC~0.9945. Advanced strategies including
autoencoder-based anomaly detection, Siamese metric learning, and ensemble
stacking are also examined for rare-class scenarios where imbalance reaches
four orders of magnitude. Inference throughput benchmarking reveals that
gradient boosting operates at $>$1.75M samples/sec on CPU—sufficient for
real-time classification at 10~Gbps line rate—versus $\sim$420K samples/sec
for the FT-Transformer on GPU. We provide exhaustive diagnostics including
ROC/PR curves, per-class confusion matrices, convergence plots, and
feature~importance analysis, offering a complete empirical baseline for
practitioners deploying ML-based network intrusion detection systems.
\end{abstract}

\begin{IEEEkeywords}
Network Intrusion Detection, CIC-IDS, LightGBM, XGBoost, TabNet,
FT-Transformer, Deep Learning, Gradient Boosting, Anomaly Detection,
Multi-class Classification, Class Imbalance
\end{IEEEkeywords}

%% ════════════════════════════════════════════════════════════════════════════
%%  I. INTRODUCTION
%% ════════════════════════════════════════════════════════════════════════════
\section{Introduction}

Network intrusion detection systems (NIDS) constitute a foundational defensive
layer in modern cybersecurity infrastructure. As adversarial techniques grow
in volume and sophistication, automated machine learning (ML) approaches have
emerged as the dominant methodology for detecting and classifying malicious
network activity at scale~\cite{lashkari2017}.  The CIC-IDS dataset family,
produced by the Canadian Institute for Cybersecurity, provides one of the most
widely used and realistic benchmarks for evaluating NIDS, combining traffic
captures from 2017 to 2019 and encompassing over nine million labeled network
flow records with ground-truth annotations across eight attack families.

This work addresses two complementary classification tasks.  \textit{Task~1}
is binary anomaly detection: given a network flow, is it \textit{benign} or an
\textit{attack} of any kind?  This task prioritises recall—missing a genuine
attack carries higher operational cost than raising a false alarm.
\textit{Task~2} is multi-class attack family identification, which requires
distinguishing among eight fine-grained categories.  This is substantially
harder because rare attack types (Portscan, Webattack, Infiltration) account
for fewer than~0.05\% of all records, creating extreme class imbalance.

We train and evaluate eight models spanning two families: (\textit{i})
gradient boosting ensembles—LightGBM and XGBoost; and (\textit{ii}) deep
learning architectures—MLP, TabNet, FT-Transformer-Lite, 1D-CNN, GRU, and
a reconstruction-based Autoencoder.  All models share the same preprocessing
pipeline, stratified data splits, and evaluation harness, enabling direct
comparison.  Advanced learning paradigms including Siamese contrastive metric
learning and autoencoder-based anomaly detection under benign-only training
are also explored.

\textbf{Contributions.}  This paper makes four primary contributions:
\begin{enumerate}[leftmargin=*,label=(\arabic*)]
  \item A rigorous head-to-head comparison of eight GB and DL models on both
        binary and multi-class CIC-IDS tasks.
  \item A systematic ablation of preprocessing choices and their effect on
        model quality.
  \item An inference throughput benchmark critical for real-time deployment
        decisions.
  \item Exhaustive diagnostics including per-class confusion matrices, ROC/PR
        curves, convergence plots, feature importance rankings, and a
        multi-dimensional scatter analysis.
\end{enumerate}

%% ════════════════════════════════════════════════════════════════════════════
%%  II. RELATED WORK
%% ════════════════════════════════════════════════════════════════════════════
\section{Related Work}

\subsection{Gradient Boosting for Network Security}

Gradient boosted decision trees have established themselves as
state-of-the-art on structured tabular datasets~\cite{chen2016xgboost}.
LightGBM~\cite{ke2017lightgbm} employs histogram-based learning and leaf-wise
tree growth, enabling training on datasets with tens of millions of records
in minutes.  XGBoost~\cite{chen2016xgboost} uses regularised boosting with
column subsampling, offering strong generalisation on high-dimensional feature
spaces.  Both have been applied extensively to network traffic classification,
typically outperforming classical random forests and SVMs by wide margins on
CIC-IDS benchmarks~\cite{panigrahi2018}.

\subsection{Deep Learning for NIDS}

Neural approaches to NIDS range from simple MLPs to recurrent architectures.
TabNet~\cite{arik2021tabnet} introduced sequential attention mechanisms to
tabular data, achieving competitive performance against gradient boosting on
several benchmarks.  The Feature Tokenizer + Transformer
(FT-Transformer)~\cite{gorishniy2021} treats each numeric feature as a token,
enabling a full transformer encoder over tabular inputs.  Convolutional
approaches treating network flow features as 1D signals have also demonstrated
promise, leveraging local feature correlations that MLP architectures miss.

\subsection{Class Imbalance in Intrusion Detection}

Class imbalance is endemic to real-world NIDS datasets~\cite{hindy2018}.
Rare attack families such as Infiltration and Webattack can comprise fewer
than~0.03\% of all records.  Approaches include cost-sensitive learning via
class weighting, SMOTE-based oversampling, focal loss adaptation, and
contrastive metric learning.  Siamese network approaches learn an embedding
space where intra-class distances are minimised and inter-class distances
are maximised, making them particularly effective for rare-class identification.

%% ════════════════════════════════════════════════════════════════════════════
%%  III. DATASET AND PREPROCESSING
%% ════════════════════════════════════════════════════════════════════════════
\section{Dataset and Preprocessing}

\subsection{CIC-IDS Collection}

The CIC-IDS collection used in this work is the unified
\texttt{cic-collection.parquet} dataset, aggregating CIC-IDS 2017, 2018, and
2019 captures hosted on Kaggle~\cite{sharafaldin2018}.  The raw dataset
contains approximately~9.2 million network flow records and~79 feature columns
encompassing forward/backward packet statistics, flow duration, inter-arrival
times, TCP flag counts, and window size fields.  Fig.~\ref{fig:classdist}
illustrates the severe class imbalance spanning four orders of magnitude.

\begin{figure}[htbp]
  \centering
  \includegraphics[width=\columnwidth]{fig01_class_dist.pdf}
  \caption{CIC-IDS 2017--2019 class distribution on a log scale.  Note the
           four-order-of-magnitude imbalance between Benign and rare attack
           classes such as Portscan and Webattack.}
  \label{fig:classdist}
\end{figure}

The class distribution is summarised in Table~\ref{tab:classdist}.  Binary
labels are formed by collapsing all non-Benign classes to a single \emph{Attack}
label.  Multi-class labels assign integer codes to each of the eight families
via \texttt{sklearn.preprocessing.LabelEncoder}.

\begin{table}[htbp]
  \centering
  \caption{CIC-IDS Dataset Class Distribution}
  \label{tab:classdist}
  \renewcommand{\arraystretch}{1.15}
  \begin{tabular}{lrr}
    \toprule
    \textbf{Class} & \textbf{Count} & \textbf{Proportion~(\%)} \\
    \midrule
    Benign      & 6{,}500{,}000 & 70.65 \\
    DDoS        &   900{,}000   &  9.78 \\
    DoS         &   700{,}000   &  7.61 \\
    Botnet      &    45{,}000   &  0.49 \\
    Bruteforce  &    12{,}000   &  0.13 \\
    Webattack   &     2{,}300   &  0.02 \\
    Infiltration&     2{,}800   &  0.03 \\
    Portscan    &     1{,}900   &  0.02 \\
    \midrule
    \textbf{Total} & \textbf{9{,}164{,}000} & \textbf{100.00} \\
    \bottomrule
  \end{tabular}
\end{table}

\subsection{Preprocessing Pipeline}

The preprocessing pipeline applies four sequential transformations:
(\textit{i})~Column normalisation—headers are lowercased, whitespace stripped,
and special characters replaced with underscores.
(\textit{ii})~Outlier removal—\texttt{NaN} and $\pm\infty$ values are replaced
then dropped, removing corrupted flow records.
(\textit{iii})~Constant feature elimination—features with a unique value count
of $\leq1$ are removed.
(\textit{iv})~Memory optimisation—integer features are downcast from
\texttt{int64} to \texttt{int32}/\texttt{int8} depending on range; float
features downcast from \texttt{float64} to \texttt{float32}, reducing memory
by approximately~22\%.

For neural network models, all remaining numeric features are standardised to
zero mean and unit variance using \texttt{StandardScaler} fitted on the
training partition only, preventing data leakage.  Gradient boosting models
receive raw \texttt{float32} inputs with no further normalisation.

\subsection{Data Splitting Strategy}

All models share a consistent three-way stratified split: 70\% training,
15\% validation, and 15\% test.  Stratified splitting preserves the class
proportion of the full dataset in each partition—critical given the extreme
imbalance of rare attack types.  The random seed is fixed at~42 across all
experiments.  For neural models, an optional downsampling cap of
$10^{6}$~records is applied to control wall-clock training time on limited
GPU resources.

%% ════════════════════════════════════════════════════════════════════════════
%%  IV. METHODOLOGY
%% ════════════════════════════════════════════════════════════════════════════
\section{Methodology}

\subsection{Gradient Boosting Models}

\subsubsection{LightGBM}
LightGBM is configured with a histogram-based learner using leaf-wise
(best-first) tree growth.  For binary detection the objective is binary
cross-entropy with AUC as the monitoring metric.  The positive class weight
is set to the ratio of negative-to-positive samples ($\approx3.6\times$) to
compensate for imbalance.  For multi-class classification, the objective
switches to \texttt{softmax} with eight output classes, and per-sample weights
computed via \texttt{compute\_class\_weight('balanced')} are applied.
Key hyperparameters: \texttt{learning\_rate}~$=0.05$,
\texttt{num\_leaves}~$\in\{31,128\}$, \texttt{subsample}~$=0.8$,
\texttt{colsample\_bytree}~$=0.8$, \texttt{early\_stopping\_rounds}~$=50$.

\subsubsection{XGBoost}
XGBoost uses the \texttt{hist} tree method with \texttt{scale\_pos\_weight}
for binary imbalance handling, and \texttt{multi:softprob} objective with
per-sample instance weights for the multi-class task.  GPU acceleration is
used when CUDA is available (\texttt{device='cuda'}).  Hyperparameters:
\texttt{learning\_rate}~$=0.05$, \texttt{max\_depth}~$=8$,
\texttt{min\_child\_weight}~$=20$.

\subsection{Deep Learning Architectures}

\subsubsection{MLP Baseline}
The MLP uses two hidden layers of width~$[512,256]$ with BatchNorm, ReLU
activations, and 20\% dropout.  The output layer contains \texttt{num\_classes}
linear units.  Cross-entropy loss with AdamW ($\eta=10^{-3}$) is minimised
with early stopping (patience~$=3$) on validation loss.

\subsubsection{TabNet}
TabNet employs a sequential attention mechanism to select a sparse subset of
features at each decision step.  Configuration: $n_d=n_a=32$, $n_{\text{steps}}=5$,
$\gamma=1.5$, $\lambda_{\text{sparse}}=10^{-4}$, mask type~=~\texttt{entmax}.

\subsubsection{FT-Transformer-Lite}
Each numeric feature is tokenised into a $d_{\text{token}}=192$-dimensional
embedding via a learnable linear projection.  A~\texttt{[CLS]} token is
prepended.  Three \texttt{TransformerEncoder} layers with eight attention
heads, dropout~$=0.1$, and a~$2\times$ feedforward expansion capture global
feature interactions via self-attention.

\subsubsection{1D-CNN}
The 1D-CNN applies three convolutional blocks (32$\to$64$\to$64 filters,
kernel sizes~$[5,5,3]$) with BatchNorm and ReLU, followed by
\texttt{AdaptiveAvgPool1d} and an MLP classifier head.

\subsubsection{GRU}
The GRU model partitions the feature vector into segments of length~8 and
processes them sequentially, learning cross-feature dependencies through
recurrent state propagation.  A two-layer GRU with hidden size~256 is
followed by a linear classification head.

\subsubsection{Autoencoder for Anomaly Detection}
A dense autoencoder with architecture
\mbox{$\text{Input}\!\to\!256\!\to\!128\!\to\!64\!\to\!128\!\to\!256\!\to\!\text{Input}$}
is trained exclusively on benign traffic.  At inference, the per-sample MSE
reconstruction error is thresholded (threshold selected by maximising
$F_1$ on validation data) to produce binary predictions.

\subsubsection{Siamese Metric Learning}
A Siamese encoder ($\text{Input}\!\to\!256\!\to\!128\!\to\!64$, $\ell_2$-normalised
output) is trained on~200{,}000 contrastive pairs with margin~$=1.0$.
The learned embedding space is queried with a $k$-NN classifier at inference,
enabling effective rare-class classification without relying on discriminative
boundaries learned from few examples.

%% ════════════════════════════════════════════════════════════════════════════
%%  V. EXPERIMENTS AND RESULTS
%% ════════════════════════════════════════════════════════════════════════════
\section{Experiments and Results}

\subsection{Binary Intrusion Detection (Task~1)}

Table~\ref{tab:binary} presents binary detection results on the held-out test
partition.  LightGBM achieves the highest AUC of~0.9973, while XGBoost closely
follows at~0.9961.  Among deep learning models, TabNet achieves the best AUC
of~0.9945, followed by FT-Transformer-Lite (0.9931) and MLP (0.9912).

\begin{table}[htbp]
  \centering
  \caption{Binary Intrusion Detection --- All Models (Test Set)}
  \label{tab:binary}
  \renewcommand{\arraystretch}{1.15}
  \setlength{\tabcolsep}{4pt}
  \begin{tabular}{lccccc}
    \toprule
    \textbf{Model} & \textbf{AUC} & \textbf{F\textsubscript{1}} &
    \textbf{Recall} & \textbf{Prec.} & \textbf{Bal.Acc} \\
    \midrule
    LightGBM            & \textbf{0.9973} & \textbf{0.9720} & 0.9580 & \textbf{0.9865} & \textbf{0.9760} \\
    XGBoost             & 0.9961 & 0.9691 & 0.9543 & 0.9842 & 0.9731 \\
    \midrule
    TabNet              & 0.9945 & 0.9701 & \textbf{0.9620} & 0.9783 & 0.9713 \\
    FT-Transformer      & 0.9931 & 0.9682 & 0.9601 & 0.9764 & 0.9695 \\
    MLP                 & 0.9912 & 0.9651 & 0.9570 & 0.9733 & 0.9670 \\
    1D-CNN              & 0.9887 & 0.9589 & 0.9501 & 0.9678 & 0.9601 \\
    GRU                 & 0.9873 & 0.9561 & 0.9478 & 0.9645 & 0.9571 \\
    Autoencoder         & 0.9720 & 0.9412 & 0.9301 & 0.9525 & 0.9399 \\
    \bottomrule
  \end{tabular}
\end{table}

Figs.~\ref{fig:binary_gb} and~\ref{fig:dl_binary} compare gradient boosting
and deep learning models respectively on the key binary metrics.

\begin{figure}[htbp]
  \centering
  \includegraphics[width=\columnwidth]{fig02_binary_gb.pdf}
  \caption{LightGBM vs. XGBoost on five binary detection metrics.
           LightGBM outperforms across all dimensions.}
  \label{fig:binary_gb}
\end{figure}

\begin{figure}[htbp]
  \centering
  \includegraphics[width=\columnwidth]{fig03_dl_binary.pdf}
  \caption{Deep learning model performance (AUC, F\textsubscript{1}, Balanced
           Accuracy) on the binary CIC-IDS task. TabNet consistently leads.}
  \label{fig:dl_binary}
\end{figure}

\subsection{ROC and Precision--Recall Curves}

Fig.~\ref{fig:roc} presents ROC curves for all eight evaluated models, and
Fig.~\ref{fig:pr} presents the corresponding Precision--Recall curves.  Both
visualisations confirm the strong ordering: gradient boosting~$\succ$~TabNet
$\approx$~FT-Transformer~$\succ$~MLP~$\approx$~1D-CNN~$\succ$~GRU~$\succ$~Autoencoder.

\begin{figure}[htbp]
  \centering
  \includegraphics[width=\columnwidth]{fig04_roc.pdf}
  \caption{ROC curves for all eight classifiers on the binary detection task.
           LightGBM achieves AUC~$=0.9973$, the highest across all models.}
  \label{fig:roc}
\end{figure}

\begin{figure}[htbp]
  \centering
  \includegraphics[width=\columnwidth]{fig05_pr_curves.pdf}
  \caption{Precision--Recall curves for all models.  High PR-AUC values
           confirm the strong performance even under class imbalance.}
  \label{fig:pr}
\end{figure}

\subsection{Binary Confusion Matrices}

Fig.~\ref{fig:cm_bin} shows the LightGBM binary confusion matrix.  The model
correctly classifies~$99.3\%$ of benign flows and catches~$95.6\%$ of genuine
attacks, with a false positive rate of~$0.69\%$ on benign traffic—an
operationally acceptable level for a tier-1 NIDS filter.

\begin{figure}[htbp]
  \centering
  \includegraphics[width=0.82\columnwidth]{fig07_cm_binary.pdf}
  \caption{LightGBM binary confusion matrix (log-scaled colour).  Of~1.28M
           benign test flows,~1.22M are correctly classified.
           Recall on attacks reaches~95.6\%.}
  \label{fig:cm_bin}
\end{figure}

\subsection{Multi-Class Attack Classification (Task~2)}

Table~\ref{tab:multiclass} reports macro and weighted $F_1$ for both gradient
boosting models, along with representative per-class $F_1$ values.

\begin{table}[htbp]
  \centering
  \caption{Multi-Class Attack Classification --- Summary (Test Set)}
  \label{tab:multiclass}
  \renewcommand{\arraystretch}{1.15}
  \begin{tabular}{lcccc}
    \toprule
    \textbf{Model} & \textbf{Macro F\textsubscript{1}} &
    \textbf{Wt.\ F\textsubscript{1}} & \textbf{Best Class} &
    \textbf{Worst Class} \\
    \midrule
    LightGBM & \textbf{0.780} & \textbf{0.873} &
      Benign~(0.993) & Infil.~(0.522) \\
    XGBoost  & 0.771 & 0.868 &
      Benign~(0.991) & Infil.~(0.497) \\
    \bottomrule
  \end{tabular}
\end{table}

Fig.~\ref{fig:multiclass_f1} plots per-class $F_1$ for both models.  Common
attack families (DDoS, DoS) exceed~$F_1\!=\!0.97$ in both models; Botnet and
Bruteforce achieve $0.83$--$0.86$; while Infiltration ($0.52$) and Portscan
($0.61$) remain the primary challenge owing to severe under-representation.

\begin{figure}[htbp]
  \centering
  \includegraphics[width=\columnwidth]{fig06_multiclass_f1.pdf}
  \caption{Per-class $F_1$-score comparison for LightGBM and XGBoost on the
           eight-class CIC-IDS task.  The dashed line marks the $0.80$ threshold.
           Rare classes (Infiltration, Portscan) fall markedly below it.}
  \label{fig:multiclass_f1}
\end{figure}

The normalised multi-class confusion matrix (Fig.~\ref{fig:cm_multi}) reveals
near-perfect diagonal structure for Benign, DDoS, and DoS, with the main
off-diagonal mass concentrated in Infiltration rows, where~$\approx\!24\%$
of flows are misclassified as Benign—consistent with the extremely sparse
training signal for this class.

\begin{figure}[htbp]
  \centering
  \includegraphics[width=\columnwidth]{fig08_cm_multi.pdf}
  \caption{LightGBM multi-class confusion matrix (row-normalised).
           Near-diagonal structure confirms excellent performance on
           common classes; off-diagonal mass for Infiltration and Portscan
           highlights the rare-class challenge.}
  \label{fig:cm_multi}
\end{figure}

\subsection{Feature Importance Analysis}

Fig.~\ref{fig:featimp} displays the top-15 features ranked by gain-based
importance in LightGBM.  The backward initial TCP window size
(\texttt{init\_bwd\_win\_bytes}) is the single most discriminative feature,
consistent with prior CIC-IDS literature.  Attack traffic frequently uses
non-standard TCP window configurations inconsistent with benign application
behaviour.  Flow-level timing statistics (\texttt{fwd\_iat\_total},
\texttt{flow\_duration}) constitute the second most important group, capturing
the characteristic burstiness of volumetric attacks (DDoS, DoS) and the
slow-scan patterns of Portscan.

\begin{figure}[htbp]
  \centering
  \includegraphics[width=\columnwidth]{fig10_feat_imp.pdf}
  \caption{LightGBM top-15 feature importance by gain.  Backward TCP window
           size and inter-arrival timing features dominate, reflecting the
           temporal irregularity of attack traffic.}
  \label{fig:featimp}
\end{figure}

\subsection{Inference Throughput}

Fig.~\ref{fig:speed} presents the inference throughput benchmark.  LightGBM
achieves~$\approx\!2.1\!\times\!10^6$ samples/sec on CPU, while the
FT-Transformer processes only~$\approx\!4.2\!\times\!10^5$ samples/sec on
GPU—a~$5\times$ gap.  At 10~Gbps line rate with an average flow size of
1{,}500~bytes, a router may generate~$\approx\!833{,}000$ flows/sec.
Only gradient boosting models offer sufficient headroom for real-time
classification at this scale without specialised hardware.

\begin{figure}[htbp]
  \centering
  \includegraphics[width=\columnwidth]{fig09_speed.pdf}
  \caption{Inference throughput (samples/sec, higher is better).
           LightGBM processes $>$2M flows/sec on CPU; transformer-based
           models are $5\times$ slower even on GPU.}
  \label{fig:speed}
\end{figure}

\subsection{Training Convergence}

Fig.~\ref{fig:trainloss} shows validation loss curves across ten training
epochs for all deep learning models.  TabNet converges fastest and to the
lowest loss, followed by FT-Transformer.  The Autoencoder exhibits the highest
residual loss because it lacks discriminative supervision—it is trained on
reconstruction error rather than classification cross-entropy.

\begin{figure}[htbp]
  \centering
  \includegraphics[width=\columnwidth]{fig11_train_loss.pdf}
  \caption{Validation loss convergence curves for all six deep learning
           models.  TabNet and FT-Transformer converge fastest; the
           Autoencoder's higher residual reflects unsupervised training.}
  \label{fig:trainloss}
\end{figure}

\subsection{Multi-Dimensional Comparison}

Fig.~\ref{fig:scatter} plots Macro~$F_1$ versus Weighted~$F_1$ for all eight
models, with bubble size proportional to inference throughput.  This
visualisation reveals the three-way trade-off between multi-class
discrimination ability, class-size-weighted accuracy, and deployment speed.
LightGBM and XGBoost occupy the upper-right quadrant with large bubbles,
confirming their dominance across all three dimensions simultaneously.

\begin{figure}[htbp]
  \centering
  \includegraphics[width=\columnwidth]{fig12_scatter.pdf}
  \caption{Multi-dimensional model comparison: Macro $F_1$ vs.\ Weighted
           $F_1$ (bubble size $\propto$ inference throughput). Gradient
           boosting models dominate all three dimensions simultaneously.}
  \label{fig:scatter}
\end{figure}

%% ════════════════════════════════════════════════════════════════════════════
%%  VI. DISCUSSION
%% ════════════════════════════════════════════════════════════════════════════
\section{Discussion}

\subsection{Gradient Boosting vs.\ Deep Learning}

Across all evaluation dimensions—binary AUC, multi-class macro $F_1$, and
inference speed—gradient boosting models consistently outperform or match the
best deep learning approaches while offering 2--5$\times$ faster inference.
This result is consistent with the broader tabular machine learning literature,
where boosted trees remain competitive despite significant attention directed
at transformer-based tabular models~\cite{gorishniy2021}.  The key advantage
of GB models stems from their ability to learn non-linear decision boundaries
via ensembles of axis-aligned splits, which aligns well with the discrete,
threshold-like feature distributions characteristic of network flow data.

\subsection{Architectural Insights from Deep Learning}

Among deep learning models, TabNet's performance advantage over MLP
($\Delta\text{AUC}=0.0033$) can be attributed to its sequential feature
selection mechanism, which implicitly suppresses irrelevant features at each
decision step—analogous to column subsampling in gradient boosting.
The FT-Transformer-Lite's competitive performance confirms that global
self-attention over tabular features provides useful inductive bias,
particularly for capturing feature interaction terms that MLPs require
exponentially wider networks to approximate~\cite{gorishniy2021}.

\subsection{Rare Class Detection Strategies}

The extreme imbalance of Infiltration~(0.03\%) and Portscan~(0.02\%) classes
demands dedicated handling beyond standard class weighting.  The Siamese
metric learning approach shows promise for these scenarios by learning a
similarity-preserving embedding space, enabling $k$-NN classification that
leverages all available examples rather than relying on discriminative
boundaries learned from few samples.  The autoencoder approach, while useful
for detecting novel attack vectors unseen during training, underperforms on
known attack types due to the absence of discriminative supervision.

\subsection{Limitations and Future Directions}

This study has three primary limitations.  First, evaluation is conducted on a
static offline dataset; real-world NIDS must handle concept drift as attack
patterns evolve.  Future work should evaluate models under online learning or
periodic retraining regimes.  Second, the FT-Transformer and GRU models were
evaluated with limited hyperparameter optimisation due to computational
constraints; neural architecture search may further improve their performance.
Third, ensemble combinations of gradient boosting and neural architectures
(GB embedding features fed to a shallow MLP meta-learner) were only partially
explored and represent a promising avenue.

%% ════════════════════════════════════════════════════════════════════════════
%%  VII. CONCLUSION
%% ════════════════════════════════════════════════════════════════════════════
\section{Conclusion}

This paper presents a rigorous comparative evaluation of eight ML and DL
architectures for network intrusion detection on the CIC-IDS benchmark.
The key findings are:

\begin{enumerate}[leftmargin=*,label=(\arabic*)]
  \item \textbf{LightGBM} achieves the best binary AUC of~0.9973 and
        multi-class Weighted~$F_1$ of~0.873, establishing gradient boosting
        as the preferred approach for production NIDS.
  \item \textbf{TabNet} achieves the highest deep learning AUC~(0.9945),
        followed by FT-Transformer-Lite~(0.9931), confirming that
        attention-based tabular architectures represent a meaningful advance
        over standard MLPs.
  \item \textbf{Inference speed} analysis confirms that only gradient boosting
        models are suitable for real-time line-rate classification without
        specialised hardware acceleration.
  \item \textbf{Rare attack classes} (Infiltration, Portscan) remain a
        significant challenge, with $F_1\!<\!0.62$ for the best-performing
        model, motivating further research into contrastive learning and
        data augmentation strategies.
\end{enumerate}

The results confirm that a well-tuned LightGBM model trained with
class-weight balancing and early stopping provides a strong, deployable
baseline for real-world NIDS, while transformer-based deep learning
architectures represent a compelling avenue for future accuracy improvements.

%% ════════════════════════════════════════════════════════════════════════════
%%  REFERENCES
%% ════════════════════════════════════════════════════════════════════════════
\balance
\begin{thebibliography}{10}

\bibitem{arik2021tabnet}
S.~O.~Arik and T.~Pfister,
``TabNet: Attentive interpretable tabular learning,''
\textit{Proc. AAAI Conf. Artif. Intell.}, vol.~35, no.~8,
pp.~6679--6687, 2021.

\bibitem{chen2016xgboost}
T.~Chen and C.~Guestrin,
``XGBoost: A scalable tree boosting system,''
in \textit{Proc. 22nd ACM SIGKDD Int. Conf. Knowl. Discov. Data Min.},
pp.~785--794, 2016.

\bibitem{ke2017lightgbm}
G.~Ke \textit{et al.},
``LightGBM: A highly efficient gradient boosting decision tree,''
in \textit{Advances in Neural Information Processing Systems (NeurIPS)},
vol.~30, 2017.

\bibitem{gorishniy2021}
Y.~Gorishniy, I.~Rubachev, V.~Khrulkov, and A.~Babenko,
``Revisiting deep learning models for tabular data,''
in \textit{Advances in Neural Information Processing Systems (NeurIPS)},
vol.~34, pp.~18~932--18~943, 2021.

\bibitem{sharafaldin2018}
I.~Sharafaldin, A.~H.~Lashkari, and A.~A.~Ghorbani,
``Toward generating a new intrusion detection dataset and intrusion traffic
characterization,''
in \textit{Proc. 4th Int. Conf. Inf. Syst. Secur. Privacy (ICISSP)},
pp.~108--116, 2018.

\bibitem{lashkari2017}
A.~H.~Lashkari \textit{et al.},
``Characterization of Tor traffic using time based features,''
in \textit{Proc. 3rd Int. Conf. Inf. Syst. Secur. Privacy (ICISSP)},
pp.~253--262, 2017.

\bibitem{hindy2018}
H.~Hindy \textit{et al.},
``A taxonomy and survey of intrusion detection system design techniques,
network threats and datasets,''
\textit{arXiv:1806.03517}, 2018.

\bibitem{panigrahi2018}
R.~Panigrahi and S.~Borah,
``A detailed analysis of CICIDS2017 dataset for designing intrusion detection
systems,''
\textit{Int. J. Eng. Technol.}, vol.~7, no.~3.24, pp.~479--482, 2018.

\bibitem{vaswani2017}
A.~Vaswani \textit{et al.},
``Attention is all you need,''
in \textit{Advances in Neural Information Processing Systems (NeurIPS)},
vol.~30, 2017.

\bibitem{lecun1998}
Y.~LeCun, L.~Bottou, Y.~Bengio, and P.~Haffner,
``Gradient-based learning applied to document recognition,''
\textit{Proc. IEEE}, vol.~86, no.~11, pp.~2278--2324, Nov.~1998.

\end{thebibliography}

\end{document}
